% ===================================================================
%  اللوحة رقم ٣ — الطفولة والشباب.
%  Traduction 2026 d'apres texto/io, controlee sur texto/fr, texto/en
%  et texto/es. Consigne : voir texto/ar/10-lawha-01.tex.
%
%  Deux scenes, et l'ido subdivise beaucoup plus que le francais :
%  huit des intertitres n'ont pas de vis-a-vis francais. L'arabe suit
%  l'ido, comme les autres traductions.
% ===================================================================

%%K t03-apar-1 apar
\VUsaut{3.0mm}
\VUtitre{15.0pt}{60}{اللوحة رقم ٣}
\VUsaut{1.6mm}
\VUtitre{11.4pt}{0}{الطفولة والشباب.}
\VUsaut{3.4mm}

%%K t03-apar-2 apar
(اللوحتان الثالثة والرابعة موضوعتان بحيث تعرضان، في أربعة
\VUgras{مشاهد عائلية}، ألفاظ \VUgras{الطقس} و\VUgras{العمر}
و\VUgras{الأسرة} و\VUgras{اللباس} و\VUgras{الصحة} الأساسية.)
\VUsaut{4.0mm}

%%K t03-tit-1 sub
\VUcentre{12.0pt}{0}{\textit{المشهد الأول.}}
\VUsaut{3.0mm}

%%K t03-tit-2 sub
\VUpk{10.2pt}{40}{تعميد في القرية.}
\VUsaut{2.4mm}

%%K t03-tit-3 sub
\VUpk{10.2pt}{40}{الربيع.}
\VUsaut{3.0mm}

%%K t03-c1-01-1 p
1. --- نحن في \VUgras{الربيع}، في شهر \VUgras{مارس} أو \VUgras{أبريل} أو
\VUgras{مايو}، في يوم من أيام \VUgras{الأسبوع}: \VUgras{الاثنين} أو
\VUgras{الثلاثاء} أو \VUgras{الأربعاء}. والساعة \VUgras{العاشرة} صباحًا.

%%K t03-c1-02-1 p
2. --- \VUgras{الطقس} جميل و\VUgras{الهواء} صافٍ. والشمس مشرقة. وبضع
\VUgras{سُحيبات} \textsuperscript{(2)} ترتفع في \VUgras{السماء} \textsuperscript{(1)}
الزرقاء.

%%K t03-c1-03-1 p
3. --- ليس على \VUgras{الأشجار} \VUgras{ورق} يُذكر. و\VUgras{الحور}
\textsuperscript{(3)} النحيل و\VUgras{الدُّلْب} \textsuperscript{(17)} الباسق لا
يحملان بعد إلا البراعم.

%%K t03-c1-04-1 p
4. --- عادت \VUgras{السنونوات} \textsuperscript{(4)}، وهي تحوم بصياح فَرِح حول
\VUgras{القمّة} \textsuperscript{(7)} التي تعلو \VUgras{برج الأجراس}
\textsuperscript{(5)} فوق \VUgras{كنيسة} \textsuperscript{(6)} القرية.

%%K t03-tit-4 sub
\VUsaut{3.4mm}
\VUpk{10.2pt}{40}{الكنيسة.}
\VUsaut{3.0mm}

%%K t03-c1-05-1 p
5. --- بسيطة جدًا هذه الكنيسة بـ\VUgras{بابها} \textsuperscript{(13)} الكبير
و\VUgras{ساعتها} \textsuperscript{(8)} الكبيرة. \VUgras{العقرب} \textsuperscript{(9)}
الصغير على الرقم X: وهو يدل على \VUgras{الساعات}. و\VUgras{الكبير}
\textsuperscript{(10)} على XII (الظهر)؛ وهو يدل على \VUgras{الدقائق}.

%%K t03-c1-06-1 p
6. --- وفي البرج تُقرَع \VUgras{الأجراس} \textsuperscript{(11)} قرعًا مَرِحًا.
وفي أعلى السطح \VUgras{صليب} \textsuperscript{(12)} حجري صغير.

%%K t03-c1-07-1 p
7. --- وللكنيسة، إلى جانب \VUgras{بابها} \textsuperscript{(13)} الكبير، باب
جانبي صغير. ومن هذا الباب يخرج \VUgras{كاهن الرعية} \textsuperscript{(15)} في
\VUgras{جُبّته} من \VUgras{غرفة الثياب} \textsuperscript{(14)} متجهًا إلى
\VUgras{دار الكاهن} \textsuperscript{(19)}.

%%K t03-c1-08-1 p
8. --- وقربه ها هي \VUgras{بائعة اللبن} \textsuperscript{(24)} مع \VUgras{حمارها}
\textsuperscript{(23)}. عائدة من المدينة حيث حملت لبنها الطيب للأطفال.

%%K t03-tit-5 sub
\VUsaut{4.0mm}
\VUpk{10.2pt}{40}{البستان.}
\VUsaut{3.4mm}

%%K t03-c1-09-1 p
9. --- والسيد الرمادي (الحمار) راضٍ عن جولة الصباح هذه. يستنشق في لذة
هواءً عطّره \VUgras{الزهر} \textsuperscript{(20)} الذي يكسو \VUgras{أشجار
الفاكهة} \textsuperscript{(18)} في \VUgras{البستان} المجاور. فتلك
\VUgras{أشجار الخوخ} \textsuperscript{(18)} و\VUgras{أشجار المشمش}
\textsuperscript{(19)} كلها وردية وبيضاء، وهي تَعِد بمحصول طيب.

%%K t03-c1-10-1 p
10. --- وفي هذه الأثناء تذهب \VUgras{النحلات} \textsuperscript{(22)} الصغيرة من
\VUgras{الخلية} \textsuperscript{(21)} إلى هناك لتجني \VUgras{عسلها} الحلو وهي
تطنّ.

%%K t03-c1-11-1 p
11. --- لكن ماذا يجري؟ ها هم أطفال القرية يُقبِلون راكضين فيبعثرون
\VUgras{الإوَزّ} \textsuperscript{(25)} المذعور. إنه الموكب يخرج من الكنيسة بعد
\VUgras{تعميد} ابن الموثّق.

%%K t03-c1-12-1 p
12. --- ويستيقظ يعقوب الصغير في \VUgras{مهده} \textsuperscript{(26)} على قرع
الأجراس الفَرِح، وتطلب أخته مادلين، وهي تريد أن ترى الموكب أيضًا، من أمها
أن تأتي فتمشّط شعرها وتلبسها على عتبة الدار.

%%K t03-c1-13-1 p
13. --- فتقبل الأم. تضع \VUgras{منديل رأسها} \textsuperscript{(28)}، وتلبس
\VUgras{صُدرتها} \textsuperscript{(29)} و\VUgras{مئزرها} \textsuperscript{(30)}،
وتجلس على كرسي صغير أمام الباب. وليس على مادلين إلا \VUgras{تنّورتها
الداخلية} \textsuperscript{(31)} و\VUgras{صُدَيريتها} \textsuperscript{(32)}.

%%K t03-c1-14-1 p
14. --- ما أسعدها! لقد نسيت أزهارها المفضلة: \VUgras{قرنفلها}
\textsuperscript{(34)} الذي تحطّ عليه \VUgras{الفراشات} \textsuperscript{(35)}،
و\VUgras{زنبقها الهولندي} \textsuperscript{(36)}، و\VUgras{زنبق الوادي}
\textsuperscript{(37)} في \VUgras{أصصه} \textsuperscript{(38)} على \VUgras{الجدار}
\textsuperscript{(33)}، و\VUgras{ورودها} \textsuperscript{(39)} التي تصنع سياجًا
بديعًا. إنها تحدّق في ثياب السيدات الفاخرة في الموكب.

%%K t03-c1-15-1 p
15. --- أما أولاد القرية فلا يعبأون كثيرًا بالثياب. يندفعون ويتزاحمون
لينالوا \VUgras{الحلوى} \textsuperscript{(55)} أو \VUgras{القطع النقدية}
\textsuperscript{(52)}. وأحدهم يبكي \textsuperscript{(40)}.

%%K t03-c1-16-1 p
16. --- وآخر داس على قائمة \VUgras{الكلب} \textsuperscript{(42)}، فانطلق يعوي
وأطار \VUgras{الدجاجة} \textsuperscript{(43)} و\VUgras{كتاكيتها} \textsuperscript{(44)}.

%%K t03-tit-6 sub
\VUsaut{5.0mm}
\VUpk{10.2pt}{40}{موكب التعميد.}
\VUsaut{4.0mm}

%%K t03-c1-17-1 p
17. --- في مقدمة الموكب تمشي \VUgras{المرضعة} \textsuperscript{(45)}. عليها
\VUgras{قُبَّعة} \textsuperscript{(46)} جميلة ذات شرائط طويلة. وهي تحمل
\VUgras{المولود} \textsuperscript{(47)} مكسوًّا كله بـ\VUgras{الدانتيل}
\textsuperscript{(48)}.

%%K t03-c1-18-1 p
18. --- وخلف المرضعة يأتي \VUgras{العرّاب} \textsuperscript{(49)}
و\VUgras{العرّابة} \textsuperscript{(53)}. وهما يوزعان الحلوى والنقود. وعلى
العرّاب \VUgras{قفازان} \textsuperscript{(51)} و\VUgras{قبعة عالية}
\textsuperscript{(50)}. وفي يد العرّابة \VUgras{باقة} \textsuperscript{(54)} جميلة.

%%K t03-c1-19-1 p
19. --- وخلفهما نرى \VUgras{عم} \textsuperscript{(56)} الطفل و\VUgras{عمته}
\textsuperscript{(58)}. على العمة \VUgras{قبعة} \textsuperscript{(59)} أنيقة، وعلى
العم \VUgras{معطف} \textsuperscript{(57)} أسود وصُدرة بيضاء. ثم يأتي
\VUgras{ابن العم} \textsuperscript{(60)} و\VUgras{ابنة العم} \textsuperscript{(61)}.
وهذه تمسك بيد \VUgras{أخت} \textsuperscript{(62)} المولود، بيرتا الصغيرة.

%%K t03-c1-20-1 p
20. --- و\VUgras{أخو} \textsuperscript{(63)} بيرتا الآخر يمشي خلفهم. يعطي يده
لـ\VUgras{قريبة} \textsuperscript{(64)}. ويأتي بعدهم \VUgras{أصدقاء}
\textsuperscript{(65)} الأسرة.

%%K t03-tit-7 sub
\VUsaut{4.6mm}
\VUpk{10.2pt}{40}{المتفرجون.}
\VUsaut{3.4mm}

%%K t03-c1-21-1 p
21. --- قرب الموكب ها هو \VUgras{شحّاذ} \textsuperscript{(66)} متكئ على
\VUgras{عكازه} \textsuperscript{(67)}. يسأل الصدقة.

%%K t03-c1-22-1 p
22. --- وعلى الجانب الآخر ولد صغير \VUgras{مهذّب} \textsuperscript{(68)} جدًا.
يرفع قبعته ويقول \VUgras{«صباح الخير»} لمن يعرفهم.

%%K t03-c1-23-1 p
23. --- وقد خرج أهل القرية من بيوتهم ليروا موكب التعميد يمر. فنرى
\VUgras{فلاحًا} \textsuperscript{(69)} في \VUgras{طاقيته القطنية} وإلى جانبه
\VUgras{فلاحة} \textsuperscript{(70)}. ثم \VUgras{عجوزًا} \textsuperscript{(71)} متكئة
على \VUgras{عصاها} \textsuperscript{(72)}. وأخيرًا \VUgras{الطحّان}
\textsuperscript{(73)} بـ\VUgras{قبعته اللبّادية} \textsuperscript{(74)} العريضة،
وفلاحة شابة في \VUgras{قباقيب} \textsuperscript{(75)} تعلّم رضيعها المشي.

%%K t03-c1-24-1 p
24. --- إنه مشهد مفعم بالحياة.
\VUsaut{4.0mm}

%%K t03-tit-8 sub
\VUcentre{12.0pt}{0}{\textit{المشهد الثاني.}}
\VUsaut{3.0mm}

%%K t03-tit-9 sub
\VUpk{10.2pt}{40}{عيد عام.}
\VUsaut{2.4mm}

%%K t03-tit-10 sub
\VUpk{10.2pt}{40}{الألعاب المائية وسباق القوارب.}
\VUsaut{3.0mm}

%%K t03-c2-01-1 p
1. --- جاء \VUgras{الصيف}. وشمس شهر \VUgras{يونيو} (يوليو أو
\VUgras{أغسطس}) المحرقة تغري الشباب بالسباحة وركوب القوارب.

%%K t03-c2-02-1 p
2. --- على الضفة اليمنى من النهر يخلع المجدّفون ثيابهم ليشاركوا في
الألعاب المائية وسباق القوارب.

%%K t03-c2-03-1 p
3. --- أحدهم يخلع \VUgras{حذاءه} \textsuperscript{(1)}؛ يحلّ رباطه (يفكّ أزراره).
وصاحباه جاهزان من قبلُ. ليس عليهما الآن إلا \VUgras{القميص الرياضي}
\textsuperscript{(2)} و\VUgras{سروال السباحة} \textsuperscript{(3)}. وهما يتحدثان في
انتظار رفاقهما. يتحدثان عن السوق والعيد المقامَين على الضفة الأخرى.

%%K t03-c2-04-1 p
4. --- وعلى الأرض قربهما \VUgras{ثياب} رفاقهما: \VUgras{زوج}
\VUgras{جزم} \textsuperscript{(4)}، و\VUgras{أحذية} \textsuperscript{(5)}،
و\VUgras{جوارب} \textsuperscript{(6)}، وقبعات \VUgras{لبّادية} \textsuperscript{(7)}
و\VUgras{قشّية} \textsuperscript{(8)}، و\VUgras{رابطة عنق} \textsuperscript{(9)}.

%%K t03-c2-05-1 p
5. --- وقد طوى أحد الشبان \VUgras{سترته} \textsuperscript{(10)} بعناية. وهو
يضعها على منشفة سيلفّ فيها جاره بعد قليل \VUgras{بذلتيهما} معًا.

%%K t03-c2-06-1 p
6. --- والولد السادس متأخر. لا تزال \VUgras{قبعته} \textsuperscript{(11)} على
رأسه. ولم يخلع \VUgras{قميصه} \textsuperscript{(12)} ولا \VUgras{ياقته}
\textsuperscript{(13)} ولا \VUgras{كمّيه} \textsuperscript{(14)}. يمسك
\VUgras{صُدرته} \textsuperscript{(17)} بيده ولا يزال في \VUgras{بنطاله}
\textsuperscript{(16)} و\VUgras{حمّالتيه} \textsuperscript{(15)}. ولن يكون جاهزًا
للسباق التالي بالتأكيد.

%%K t03-c2-07-1 p
7. --- وقرب الشبان ينزل درب يكتنفه \VUgras{الشوك} \textsuperscript{(18)} إلى حافة
الماء، حيث يهيئ يعقوب العجوز، بين \VUgras{القصب} \textsuperscript{(19)}،
\VUgras{الألعاب النارية} \textsuperscript{(20)} التي ستُطلَق في المساء.

%%K t03-tit-11 sub
\VUsaut{4.4mm}
\VUpk{10.2pt}{40}{السباق.}
\VUsaut{3.4mm}

%%K t03-c2-08-1 p
8. --- وعلى النهر تتنزه بضع سيدات مستظلات بمظلاتهن في \VUgras{زورق}
\textsuperscript{(21)}. وهن يتابعن السباق، وهو مثير جدًا. ففي كل \VUgras{قارب
سباق} \textsuperscript{(22)} أربعة \VUgras{مجدّفين} \textsuperscript{(24)} يجدّفون بكل
قواهم، بينما يمسك \VUgras{الدفّاف} \textsuperscript{(26)} \VUgras{الدفّة}
\textsuperscript{(25)} ويوجّه المركب الهش. و\VUgras{المجاديف} \textsuperscript{(23)}
تضرب الماء على إيقاع واحد، والقوارب الخفيفة تنطلق بسرعة تدوخ.

%%K t03-c2-09-1 p
9. --- وبضعة شبان يتنافسون، قرب دعامة الجسر، على جائزة \VUgras{السارية
المدهونة} \textsuperscript{(28)}. أحدهم يتقدم بحذر على العمود المرن الزلق ليبلغ
\VUgras{العَلَم} \textsuperscript{(29)} الصغير المثبّت في طرفه. أما
\VUgras{منافسه} \textsuperscript{(30)} التعِس فقد سقط في الماء. وهو يسبح عائدًا
إلى الشاطئ.

%%K t03-c2-10-1 p
10. --- والفتيان الأكبر \VUgras{يتبارزون على الماء} \textsuperscript{(32)} قرب
\VUgras{الرصيف} \textsuperscript{(33)}. و\VUgras{حشد} \textsuperscript{(31)} كبير
يتفرج على هذه الألعاب كلها، متكئًا على \VUgras{سور} \VUgras{الجسر}
\textsuperscript{(27)} ومتراصًّا على طول \VUgras{الضفة}. غير أن بعض المتفرجين
يلتفتون ليتابعوا لعبة \VUgras{العمود المتسلَّق} \textsuperscript{(45)}، أو
ليعجبوا بـ\VUgras{موكب البلدية} الذاهب إلى \VUgras{توزيع الجوائز}.

%%K t03-c2-11-1 p
11. --- ففي المقدمة بضعة \VUgras{صبية} \textsuperscript{(35)} يركضون أمام
\VUgras{عازفي الفرقة} \textsuperscript{(36)}؛ ثم يأتي \VUgras{العمدة}
\textsuperscript{(37)} ونوابه و\VUgras{المجلس البلدي} للمدينة الصغيرة. وفي
الآخِر يسير \VUgras{رجال الإطفاء} \textsuperscript{(38)}.

%%K t03-tit-12 sub
\VUsaut{5.0mm}
\VUpk{10.2pt}{40}{السوق.}
\VUsaut{3.6mm}

%%K t03-c2-12-1 p
12. --- في \VUgras{ساحة السوق} \textsuperscript{(39)} زحام شديد. يأتي الناس
ليروا \VUgras{الأكشاك} المختلفة: \VUgras{الخيل الخشبية} \textsuperscript{(40)}،
و\VUgras{السيرك} \textsuperscript{(41)} حيث بدأ العرض من قبلُ؛ و\VUgras{حلبة
الرقص} \textsuperscript{(42)} حيث ينعم شبان المدينة وفتياتها بمتعة
\VUgras{الرقص}.

%%K t03-c2-13-1 p
13. --- وفي العمق نرى \VUgras{الحلبة} \textsuperscript{(44)} حيث يصارع
\VUgras{المصارع} الإسباني الشجاع \VUgras{الثور} \textsuperscript{(45)} الهائج؛
ثم \VUgras{الأفعوانية} \textsuperscript{(49)} و\VUgras{ميدان الرماية}
\textsuperscript{(46)} حيث يتمرن الناس على استعمال \VUgras{الأسلحة النارية}
والرمي على \VUgras{الهدف}.

%%K t03-c2-14-1 p
14. --- وعلى مقربة ها هو \VUgras{مسرح السوق} \textsuperscript{(47)} حيث تُمثَّل
\VUgras{الملهاة} و\VUgras{الدراما}. وقريبًا منه جدًا يقوم كشك
\VUgras{العملاق} \textsuperscript{(48)} و\VUgras{القزم}. \VUgras{طول} الأول
متران وخمسة عشر سنتيمترًا؛ والثاني بالكاد في طول طفل.

%%K t03-c2-15-1 p
15. --- وأخيرًا ها هي \VUgras{حديقة الحيوان} \textsuperscript{(50)} الكبيرة
بـ\VUgras{وحوشها}: \VUgras{نمرها} \textsuperscript{(51)} ذو \VUgras{المخالب}
القوية، و\VUgras{أسدها} \textsuperscript{(52)} ذو \VUgras{اللبدة} الكثيفة،
و\VUgras{زرافتاها} \textsuperscript{(53)}، و\VUgras{فيلها} \textsuperscript{(54)} الذي
يستعمل \VUgras{خرطومه} في براعة.

%%K t03-c2-16-1 p
16. --- و\VUgras{المروّض} \textsuperscript{(55)} الذي يدرّب هذه الحيوانات واقف
عند باب الكشك.
\VUsaut{6.0mm}
\VUfilet{22mm}
